
\item Let
$$
A=
\begin{bmatrix}
1&2&2&-1\\
2&4&-2&2\\
0&0&6&-4
\end{bmatrix},
\qquad
B=
\begin{bmatrix}
-1\\4\\-6
\end{bmatrix}.
$$

Consider the system $(S)$ of linea equations in the unknowns
$x_1,x_2,x_3,x_4$, given by $AX=B$.

\begin{enumerate}[label=\alph*)]

\item Without solving $(S)$, show that$(2,0,-3,-3)$ is a solution and that
$(2,0,0,-3)$ is not a solution.

\begin{answer}
    Let
    $$
    u=(2,0,-3,-3).
    $$

    We verify that
    $$
    \begin{aligned}
    Au
    &=
    \begin{bmatrix}
    1&2&2&-1\\
    2&4&-2&2\\
    0&0&6&-4
    \end{bmatrix}
    \begin{bmatrix}
    2\\0\\-3\\-3
    \end{bmatrix}\\[6pt]
    &=
    \begin{bmatrix}
    2-6+3\\
    4+6-6\\
    -18+12
    \end{bmatrix}\\[6pt]
    &=
    \begin{bmatrix}
    -1\\4\\-6
    \end{bmatrix}
    =B.
    \end{aligned}
    $$

    Therefore, $u$ is a solution of $(S)$.

    Now let
    $$
    v=(2,0,0,-3).
    $$

    Then
    $$
    \begin{aligned}
    Av
    &=
    \begin{bmatrix}
    1&2&2&-1\\
    2&4&-2&2\\
    0&0&6&-4
    \end{bmatrix}
    \begin{bmatrix}
    2\\0\\0\\-3
    \end{bmatrix}\\[6pt]
    &=
    \begin{bmatrix}
    5\\-2\\12
    \end{bmatrix}
    \neq B.
    \end{aligned}
    $$

    Hence $v$ is not a solution of $(S)$.

    {\color{red}
    \textbf{Note:} To verify a particular solution, it is sufficient to check $AX=B$.
    The entire system need not be solved.}
\end{answer}

\item Solve the associated homogeneous
system using Gauss-Jordan elimination.

\begin{answer}
    The associated homogeneous system is $AX=0$.

    Its augmented matrix is
    $$
    \left[
    \begin{array}{cccc|c}
    1&2&2&-1&0\\
    2&4&-2&2&0\\
    0&0&6&-4&0
    \end{array}
    \right].
    $$

    $$
    \xrightarrow{R_2\leftarrow R_2-2R_1}
    \left[
    \begin{array}{cccc|c}
    1&2&2&-1&0\\
    0&0&-6&4&0\\
    0&0&6&-4&0
    \end{array}
    \right]
    $$

    $$
    \xrightarrow{R_3\leftarrow R_3+R_2}
    \left[
    \begin{array}{cccc|c}
    1&2&2&-1&0\\
    0&0&-6&4&0\\
    0&0&0&0&0
    \end{array}
    \right]
    $$

    $$
    \xrightarrow{R_2\leftarrow-\frac16R_2}
    \left[
    \begin{array}{cccc|c}
    1&2&2&-1&0\\
    0&0&1&-\frac23&0\\
    0&0&0&0&0
    \end{array}
    \right]
    $$

    $$
    \xrightarrow{R_1\leftarrow R_1-2R_2}
    \left[
    \begin{array}{cccc|c}
    1&2&0&\frac13&0\\
    0&0&1&-\frac23&0\\
    0&0&0&0&0
    \end{array}
    \right].
    $$

    Therefore,
    $$
    r(A)=2<4=n.
    $$

    The homogeneous system is consistent and indeterminate, with degree
    $$
    p=4-2=2.
    $$

    The corresponding equations are
    $$
    \begin{cases}
    x_1+2x_2+\frac13x_4=0\\
    x_3-\frac23x_4=0
    \end{cases}
    $$

    Thus,
    $$
    \begin{cases}
    x_1=-2x_2-\frac13x_4\\
    x_3=\frac23x_4
    \end{cases}
    $$

    {\color{red}
    \textbf{Note:} The pivot columns are $1$ and $3$.
    Therefore $x_2$ and $x_4$ are free variables.}

    Let
    $$
    x_2=\alpha,\qquad
    x_4=\beta,
    \qquad\alpha,\beta\in\mathbb R.
    $$

    Then
    $$
    \begin{cases}
    x_1=-2\alpha-\frac13\beta\\
    x_2=\alpha\\
    x_3=\frac23\beta\\
    x_4=\beta
    \end{cases}
    $$

    Hence
    $$
    \begin{aligned}
    (x_1,x_2,x_3,x_4)
    ={}&
    \left(-2\alpha-\frac13\beta,
    \alpha,\frac23\beta,\beta\right)\\
    ={}&
    (-2\alpha,\alpha,0,0)\\
    &+\left(-\frac13\beta,0,
    \frac23\beta,\beta\right)\\
    ={}&
    \alpha(-2,1,0,0)\\
    &+\beta\left(-\frac13,0,
    \frac23,1\right).
    \end{aligned}
    $$

    \textbf{General solution of the associated homogeneous system}

    $$
    \boxed{
    \begin{aligned}
    (x_1,x_2,x_3,x_4)
    ={}&\alpha(-2,1,0,0)\\
    &+\beta\left(-\frac13,0,
    \frac23,1\right)
    \end{aligned}
    }
    $$
    $$
    \alpha,\beta\in\mathbb R.
    $$
\end{answer}

\item Give the general solution of the system $(S)$.

\begin{answer}
    From part a), we know that
    $$
    u=(2,0,-3,-3)
    $$
    is a particular solution of $AX=B$.

    From part b), the general solution of the associated homogeneous system is
    $$
    \begin{aligned}
    (x_1,x_2,x_3,x_4)
    ={}&\alpha(-2,1,0,0)\\
    &+\beta\left(-\frac13,0,
    \frac23,1\right).
    \end{aligned}
    $$

    Therefore, the general solution of $(S)$ is
    $$
    \boxed{
    \begin{aligned}
    (x_1,x_2,x_3,x_4)
    ={}&(2,0,-3,-3)\\
    &+\alpha(-2,1,0,0)\\
    &+\beta\left(-\frac13,0,
    \frac23,1\right)
    \end{aligned}
    }
    $$
    $$
    \alpha,\beta\in\mathbb R.
    $$

    Since the homogeneous system has two free variables, $(S)$ is consistent and
    indeterminate, with degree $2$.

    {\color{red}
    \textbf{Note:} The general solution of a consistent system is the sum of one particular solution
    and the general solution of its associated homogeneous system.}
\end{answer}

\item Give two solutions of $(S)$
different from the one considered
in part a).

\begin{answer}
    Using the general solution obtained in part c), we assign values to the free parameters.

    For $\alpha=1$ and $\beta=0$:
    $$
    \begin{aligned}
    (x_1,x_2,x_3,x_4)
    &=(2,0,-3,-3)+(-2,1,0,0)\\
    &=\boxed{(0,1,-3,-3)}.
    \end{aligned}
    $$

    For $\alpha=0$ and $\beta=3$:
    $$
    \begin{aligned}
    (x_1,x_2,x_3,x_4)
    &=(2,0,-3,-3)+(-1,0,2,3)\\
    &=\boxed{(1,0,-1,0)}.
    \end{aligned}
    $$

    Both are solutions of $(S)$ and are distinct from
    $(2,0,-3,-3)$.
\end{answer}

\end{enumerate}
